Large language models (LLMs) now generate research ideas that human experts judge more novel than their own, yet the field distrusts these ideas because ``novelty'' is poorly defined. A natural fix is to \emph{ground} the abstract notion of novelty in a concrete reasoning process. Six such mechanisms have been proposed in isolation---deduction, induction, proposal writing, outcome imagination, small experiments, and literature comparison---but no study compares them under one controlled protocol. We run the first head-to-head ablation: one fixed backbone (\gptfourone), the canonical 7 NLP ideation topics, an identical output schema across all conditions, and 5 replicates per cell, for 245 ideas. We evaluate with three independent lenses: an objective embedding metric (distance to 1{,}050 real prior-work abstracts), a 2-judge off-family LLM rubric, and position-controlled pairwise comparisons. Grounding does \emph{not} uniformly improve quality; instead each mechanism steers the novelty--feasibility trade-off in a specific way. Crucially, the apparent novelty gains exist almost entirely in the LLM judge's eye: they vanish under the objective metric, with which judge novelty is \emph{completely uncorrelated} (Spearman $\rho = 0.01$). Literature comparison wins 79\% of head-to-head comparisons and is the only mechanism that raises idea-set diversity, while ``small experiment'' grounding actively hurts novelty by pulling ideas toward already-published methods. Our results give a controlled, quantitative confirmation of the ``novelty mirage'' and argue that single-number LLM-judged novelty should never be reported alone.
